\section{Related Work}
\label{sec:related_work}

Our work is situated at the intersection of model merging, dynamic neural network ensembling, parameter-efficient fine-tuning, and biologically-inspired deep learning architectures. In this section, we review the existing literature across these domains and contrast them with EpiMerge.

\subsection{Static Model Merging}
Model merging focuses on synthesizing the parameters of multiple task-specific expert neural networks into a single cohesive model without undergoing expensive joint training or fine-tuning \cite{langley00}. A simple but surprisingly effective approach is Task Arithmetic \cite{task_arithmetic}, which constructs task-specific parameter vectors (expert weight minus base pre-trained weight) and adds their linear combination to the base pre-trained model. Other methods, such as RegMean \cite{regmean}, compute layer-wise analytical solutions using activation statistics to align weights. To mitigate the interference caused by opposing sign directions and low-magnitude values, TIES-Merging \cite{ties} introduces sign agreement resolution and magnitude thresholding. Similarly, DARE \cite{dare} applies random pruning to task vectors to sparsify them before merging. 

Additional static merging techniques resolve permutation symmetries before averaging. For example, Git Re-Basin \cite{git_rebasin} uses matching algorithms (such as the Hungarian algorithm) to align the permutable weights of independently trained models. ZipIt! \cite{zipit} goes further by matching and merging features within the same layers, allowing multi-task parameters to be merged even when model topologies differ. While these static merging methods are zero-overhead, they all force a single set of parameters for all test-time inputs. This inevitably leads to parameter conflicts and representation bottlenecks when combining experts trained on divergent distributions \cite{representation_collapse, adamerging}. EpiMerge avoids this zero-sum compromise entirely by dynamically reconstructing sample-specific weight matrices at test-time.

\subsection{Dynamic Model Merging \& Test-Time Adaptation}
To address the limitations of static merging, dynamic merging methods modify the merging parameters at inference time. A notable example is AdaMerging \cite{adamerging}, which introduces a test-time adaptation (TTA) strategy. AdaMerging optimizes layer-wise merging coefficients on-the-fly by minimizing the prediction entropy on unsupervised test batches. Although AdaMerging performs well under steady, predictable distribution streams, it is highly sensitive to test stream characteristics. As shown in recent studies on test-time adaptation vulnerabilities \cite{wang2020tent, tta_rugged, tta_shift}, unsupervised objectives can collapse under temporal non-I.I.D. shifts (e.g., sudden bursts of a single task) or extreme small-batch noise, causing severe transductive overfitting and representation collapse.

Other dynamic ensembling methods use gating networks to predict ensembling coefficients for different experts. For instance, QWS-Merge \cite{qws_merge} uses quantum-inspired wavefunction superposition to generate sample-wise merging coefficients. However, to execute the forward pass efficiently without sequentially looping through the batch, these routers average the computed coefficients across the batch dimension to produce a single weight matrix for the entire batch. This batch-averaged shortcut induces transductive batch coupling when the input batch contains samples from multiple different tasks, reducing the dynamic weights to a static average and violating sample independence. In contrast, EpiMerge performs \emph{true} sample-wise dynamic merging. Through vectorized tensor contractions, we execute a parallelized forward pass where each sample in a batch is processed by its own dynamically modulated parameter state, ensuring complete sample independence.

From a broader perspective, the concept of sample-dependent weight modulation dates back to Hypernetworks \cite{ha2016hypernetworks} and Dynamic Filter Networks \cite{jia2016dynamic}, which dynamically generate neural network parameters on-the-fly. While standard Hypernetworks require a large external network to predict high-dimensional parameters (causing extreme computational and parameter overhead), EpiMerge focuses on coordinate-wise gating of pre-trained expert vectors using a low-rank dual-masking framework ($G = \mathbf{r} \otimes \mathbf{c}$), achieving extreme parameter and calibration efficiency during multi-task synthesis.

\subsection{Mixture-of-Experts (MoE) \& Parameter-Efficient Fine-Tuning (PEFT)}
Our work also relates to Mixture-of-Experts (MoE) architectures and Parameter-Efficient Fine-Tuning (PEFT). 
Classical MoE models \cite{shazeer2017outrageously, fedus2022switch, gshard} route individual tokens or sequences to discrete, specialized feed-forward subnetworks. While highly expressive, MoE models require massive parameter footprints, specialized distributed systems, and suffer from high routing latency. Moreover, converting pre-trained dense experts into sparse MoE models requires non-trivial initialization and calibration \cite{moe_merging, merge_moe_layers}.

On the other hand, PEFT techniques such as LoRA \cite{lora}, AdaLoRA \cite{adalora}, and DoRA \cite{dora} inject low-rank trainable weight matrices into pre-trained architectures to facilitate efficient adaptation. While LoRA has been used to merge specialized adapters statically \cite{lora_hub, zip_loras}, these methods are designed for parameter-efficient \emph{training} and still lead to static merged configurations during deployment. EpiMerge utilizes the mathematical formulation of low-rank updates (i.e., row-wise and column-wise outer products) but repurposes them as a test-time dynamic ensembling module. Instead of training static low-rank adapters, we learn dynamic, sample-dependent row and column epigenetic masks that modulate pre-trained experts on-the-fly.

\subsection{Biologically-Inspired Deep Learning}
Drawing inspiration from biological systems has a long history in artificial intelligence, from the original McCulloch-Pitts neuron \cite{mcculloch_pitts} to synaptic consolidation algorithms like Elastic Weight Consolidation (EWC) \cite{ewc} designed to mitigate catastrophic forgetting.
In molecular biology, the static genome (primary DNA sequence) of an organism remains identical across all cells. However, cell differentiation and rapid environmental adaptation are controlled by the \emph{epigenome} \cite{epigenetics_textbook, chemical_genetics}. Epigenetic mechanisms, such as DNA methylation and histone acetylation, act as reversible regulatory markers that scale, silence, or amplify gene expression in response to cellular and environmental stimuli.

Despite the power of epigenetics in biology, it has received limited attention in deep learning. A few works explore epigenetic algorithms in genetic programming or neural architecture search \cite{epigenetic_gp}, but they do not apply them to the parameter spaces of deep networks. EpiMerge is the first to directly map this molecular regulatory layer to model merging. By treating the pre-trained base model as the static genetic code and the expert task vectors as specialized genes, we learn an input-dependent, reversible epigenetic masking layer (ERHs) that scales expert parameter expression on-the-fly. This delivers a highly adaptive, biologically-mimetic deep learning system.
